\documentclass[12pt, a4paper]{letter}
\usepackage[utf8]{inputenc}
\usepackage[T1]{fontenc}
\usepackage{geometry}
\usepackage{amssymb}
\usepackage{hyperref}

\geometry{margin=1in}

\hypersetup{colorlinks=true, linkcolor=blue, urlcolor=blue}

\signature{Rafael D. De Paz \\ Sovereign Architect \\ \texttt{me@rdepaz.com}}
\address{Rafael D. De Paz \\ Sovereign Architect}

\begin{document}
\begin{letter}{Editorial Board}

\opening{Dear Editors,}

I am writing to submit the enclosed manuscript, entitled

\begin{center}
\textbf{``The Grace Protocol: Dual-Track Input Validation for Graceful Cyber-Security''}
\end{center}

\noindent for consideration for publication.

\medskip

\textbf{Summary.} In formal Language-Theoretic Security (LangSec), deterministic parsers enforce absolute security but often suffer from false positives on edge-cases. This paper introduces the Grace Protocol, a dual-track validation architecture. Track A (The Justiciary) applies strict, $O(1)$ determinism to reject malformed packets. Track B (The Intercessor) utilizes heuristic, Bayesian models to analyze non-malicious drift. We prove that routing edge cases through heuristic mercy rather than immediate terminal rejection creates a computationally superior Intrusion Detection System.

\medskip

This manuscript has not been previously published and is not under consideration
at any other journal. I confirm no conflicts of interest.

\textbf{Cryptographic Lineage \& Validation.} This mathematical substrate has been cryptographically sealed and tracked on the global Sovereign Master Ledger to prevent retroactive editing and to verify the source authorship. The immutable SHA-256 integrity checksums, formal PDF renderings, and lineage authorities can be verified at \url{https://rdepaz.com/research}.

\medskip

\closing{Respectfully submitted,}
\end{letter}
\end{document}
